% ID: SEL-D0066-OB-E03
% ESTADO: cotejada. Segunda pasada visual y borrador compilado revisados.
% Fuente: Selectividad/sel_2010_matematicas_aplicadas/mat_apli/mat_apli_sept/mat_apli_sept_rb/EXAMEN-ANDALUCÍA-2.pdf
% SHA-256: 24c494a5d013dbc7f5adaac9896dee2d79f22540518359c67eac29945a26d4f0
% Página: 2; opción B; ejercicio original 3.

En un centro de enseñanza secundaria se sabe que el $45\%$ de los alumnos juegan al fútbol, que el $60\%$ practican atletismo, y que de los que practican atletismo el $50\%$ juegan al fútbol.
\begin{enumerate}[label=\alph*)]
\item \textbf{(0.75 puntos)} \textquestiondown{}Qué porcentaje de alumnos practican ambos deportes?
\item \textbf{(0.75 puntos)} Si se elige al azar un alumno de ese centro, \textquestiondown{}cuál es la probabilidad de que no practique ninguno de estos deportes?
\item \textbf{(1 punto)} Si un alumno de ese centro no juega al fútbol, \textquestiondown{}cuál es la probabilidad de que practique atletismo?
\end{enumerate}
